% =========================================================
% Induced Orders and Order Embeddings
% =========================================================

\subsection{Induced Orders and Order Embeddings}

% ---------------------------------------------------------
% The construction
% ---------------------------------------------------------

\begin{exposition}
The simplest way to put an order on a set $A$ is to compare elements of $A$
\emph{indirectly} through a function into an already ordered set. This
construction, called the \emph{induced order} or \emph{pullback order},
explains how an ambient order restricts to subsets, how functions on a common
domain acquire a pointwise order, and how order information can be transported
through maps.
\end{exposition}

\begin{definitionbox}{Definition (Induced Order)}
\begin{definition}[Induced Order]\label{def:induced-order}
Let $(B, \leq')$ be a partially ordered set and let $f : A \to B$ be a
function. The \emph{order induced by $f$} (or \emph{pullback order along
$f$}) is the relation $\leq_f$ on $A$ defined by
\[
x \;\leq_f\; y
\;\;\Longleftrightarrow\;\;
f(x) \;\leq'\; f(y).
\]
\end{definition}
\end{definitionbox}

\begin{remark*}[Standard quantified statement]
$\forall x,y\in A,\; x\leq_f y \leftrightarrow f(x)\leq' f(y)$.
\end{remark*}

\begin{remark*}[Predicate reading]
$\operatorname{InducedOrder}(\leq_f,f,\mathsf{PartiallyOrderedSet}(B,\leq'))$
means that $\leq_f$ compares elements of $A$ by applying $f$ and comparing
their images in the ordered set $(B,\leq')$.
\end{remark*}

\begin{remark*}[Interpretation]
The order on $A$ is inherited from $B$ through the map $f$.
\end{remark*}

\begin{remark*}[Exposition]
Two elements of $A$ are compared by sending them through $f$ and comparing
their images in $(B, \leq')$. The induced order contains no independent
comparison rule on $A$; every comparison is read through the ordered codomain.
\end{remark*}

\begin{dependencies}
\begin{itemize}
  \item \hyperref[def:partial-order]{Partial Order}
  \item \hyperref[def:function]{Function}
  \item \hyperref[def:relation]{Relation}
\end{itemize}
\end{dependencies}

% ---------------------------------------------------------
% Reflexivity and transitivity: always free
% ---------------------------------------------------------

\begin{propositionbox}{Proposition}
\begin{proposition}[\texorpdfstring{\hyperref[prf:induced-preorder]{$\leq_f$ is always a preorder}}{$\leq_f$ is always a preorder}]\label{prop:induced-preorder}
For any function $f : A \to B$ and partial order $(B, \leq')$, the relation
$\leq_f$ is a preorder on $A$: it is reflexive and transitive.
\hyperref[prf:induced-preorder]{\textit{Go to proof.}}
\end{proposition}
\end{propositionbox}

\begin{remark*}[Standard quantified statement]
$\forall f:A\to B,\; \mathsf{PartiallyOrderedSet}(B,\leq') \to
\operatorname{Preorder}(\leq_f,A)$.
\end{remark*}

\begin{remark*}[Predicate reading]
$\operatorname{Preorder}(\leq_f,A)$ means that the induced relation is
reflexive and transitive on $A$.
\end{remark*}

\begin{remark*}[Interpretation]
Pulling back a partial order along any function always preserves reflexivity
and transitivity.
\end{remark*}

\begin{remark*}[Exposition]
Both properties are inherited directly from $(B,\leq')$. Reflexivity at $x$
uses $f(x)\leq' f(x)$, and transitivity along $x,y,z$ uses transitivity along
$f(x),f(y),f(z)$.
\end{remark*}

\begin{dependencies}
\begin{itemize}
  \item \hyperref[def:induced-order]{Induced Order}
  \item \hyperref[def:preorder]{Preorder}
  \item \hyperref[def:function]{Function}
  \item \hyperref[def:partial-order]{Partial Order}
\end{itemize}
\end{dependencies}

% ---------------------------------------------------------
% Antisymmetry: requires injectivity
% ---------------------------------------------------------

\begin{definitionbox}{Definition ($f$-Indistinguishable Elements)}
\begin{definition}[$f$-Indistinguishable Elements]\label{def:f-indist}
Let $f : A \to B$ and $\leq_f$ be the induced order. Two elements
$x, y \in A$ are \emph{$f$-indistinguishable}, written $x \sim_f y$,
if both $x \leq_f y$ and $y \leq_f x$, equivalently if
\[
f(x) \leq' f(y) \quad \text{and} \quad f(y) \leq' f(x).
\]
\end{definition}
\end{definitionbox}

\begin{remark*}[Standard quantified statement]
$x\sim_f y \leftrightarrow (x\leq_f y \wedge y\leq_f x)
\leftrightarrow (f(x)\leq' f(y)\wedge f(y)\leq' f(x))$.
\end{remark*}

\begin{remark*}[Predicate reading]
$\operatorname{FunctionIndistinguishable}(x,y,f,\mathsf{PartiallyOrderedSet}(B,\leq'))$
means that $x$ and $y$ cannot be separated by the order after applying $f$.
\end{remark*}

\begin{remark*}[Interpretation]
$f$-indistinguishable elements occupy the same ordered position after being
viewed through $f$.
\end{remark*}

\begin{remark*}[Exposition]
When $\leq'$ is a partial order, $x \sim_f y$ forces $f(x)=f(y)$ by
antisymmetry in $B$. Thus non-injectivity is exactly the obstruction that can
make the induced preorder fail to be antisymmetric.
\end{remark*}

\begin{dependencies}
\begin{itemize}
  \item \hyperref[def:induced-order]{Induced Order}
  \item \hyperref[def:function]{Function}
  \item \hyperref[def:partial-order]{Partial Order}
\end{itemize}
\end{dependencies}

\begin{propositionbox}{Proposition}
\begin{proposition}[\texorpdfstring{\hyperref[prf:induced-poset]{$\leq_f$ is a partial order iff $f$ is injective}}{$\leq_f$ is a partial order iff $f$ is injective}]\label{prop:induced-poset}
Let $(B, \leq')$ be a partially ordered set and $f : A \to B$.
The induced order $\leq_f$ is a partial order on $A$ if and only if
$f$ is injective.
\hyperref[prf:induced-poset]{\textit{Go to proof.}}
\end{proposition}
\end{propositionbox}

\begin{remark*}[Standard quantified statement]
$\operatorname{PartialOrder}(\leq_f,A)\leftrightarrow
\operatorname{Injective}(f,A,B)$, assuming
$\mathsf{PartiallyOrderedSet}(B,\leq')$ and
$\forall x,y\in A,\; x\leq_f y\leftrightarrow f(x)\leq' f(y)$.
\end{remark*}

\begin{remark*}[Predicate reading]
$\operatorname{PartialOrder}(\leq_f,A)$ holds exactly when
$\operatorname{Injective}(f,A,B)$ holds, because antisymmetry of the pullback
is equivalent to $f$ separating distinct elements of $A$.
\end{remark*}

\begin{remark*}[Interpretation]
The induced preorder becomes a genuine partial order precisely when $f$ has no
collisions.
\end{remark*}

\begin{remark*}[Exposition]
Reflexivity and transitivity lift freely from $B$. Antisymmetry is different:
if $f$ collapses distinct points $x\neq y$, then $x\leq_f y$ and $y\leq_f x$
can both hold while $x\neq y$.
\end{remark*}

\begin{remark*}[Exposition]
Let $B=(\mathbb{N},\leq)$ and let $f:\{a,b\}\to\mathbb{N}$ be constant with
$f(a)=f(b)=0$. Then $a\leq_f b$ and $b\leq_f a$, but $a\neq b$, so
$\leq_f$ is not antisymmetric.

Let $B=(\mathcal{P}(\{1,2,3\}),\subseteq)$ and let $A=\{x,y,z\}$ with
$f(x)=\{1\}$, $f(y)=\{2\}$, and $f(z)=\{1,2\}$. Then $f$ is injective, and
the induced relation is a partial order with $x\leq_f z$, $y\leq_f z$, and
$x,y$ incomparable.
\end{remark*}

\begin{dependencies}
\begin{itemize}
  \item \hyperref[def:induced-order]{Induced Order}
  \item \hyperref[def:partial-order]{Partial Order}
  \item \hyperref[def:injective]{Injective Function}
\end{itemize}
\end{dependencies}

% ---------------------------------------------------------
% Suborders
% ---------------------------------------------------------

\begin{definitionbox}{Definition (Suborder / Restriction)}
\begin{definition}[Suborder / Restriction]\label{def:suborder}
Let $(B, \leq')$ be a partially ordered set and $S \subseteq B$. The
\emph{suborder} on $S$ (or \emph{induced order on $S$}) is the relation
$\leq'_S$ defined by
\[
x \;\leq'_S\; y
\;\;\Longleftrightarrow\;\;
x \leq' y,
\quad \text{for } x, y \in S.
\]
\end{definition}
\end{definitionbox}

\begin{remark*}[Standard quantified statement]
$\forall x,y\in S,\; x\leq'_S y \leftrightarrow x\leq' y$.
\end{remark*}

\begin{remark*}[Predicate reading]
$\operatorname{Suborder}(\leq'_S,S,\mathsf{PartiallyOrderedSet}(B,\leq'))$
means that $S$ carries the order inherited from the ambient poset $B$.
\end{remark*}

\begin{remark*}[Interpretation]
A suborder restricts the ambient comparison relation to a subset.
\end{remark*}

\begin{remark*}[Exposition]
The suborder on $S\subseteq B$ is the order induced by the inclusion map
$\iota:S\hookrightarrow B$, $\iota(x)=x$. Since $\iota$ is injective, the
suborder is always a partial order whenever $(B,\leq')$ is.
\end{remark*}

\begin{dependencies}
\begin{itemize}
  \item \hyperref[def:partial-order]{Partial Order}
  \item \hyperref[def:subset]{Subset}
  \item \hyperref[def:relation]{Relation}
\end{itemize}
\end{dependencies}

% ---------------------------------------------------------
% Order embeddings
% ---------------------------------------------------------

\begin{definitionbox}{Definition (Order Embedding)}
\begin{definition}[Order Embedding]\label{def:order-embedding}
Let $(A, \leq)$ and $(B, \leq')$ be partially ordered sets. A function
$f : A \to B$ is an \emph{order embedding} if for all $x, y \in A$,
\[
x \;\leq\; y
\;\;\Longleftrightarrow\;\;
f(x) \;\leq'\; f(y).
\]
\end{definition}
\end{definitionbox}

\begin{remark*}[Standard quantified statement]
$\operatorname{OrderEmbedding}(f,\mathsf{PartiallyOrderedSet}(A,\leq),
\mathsf{PartiallyOrderedSet}(B,\leq'))\leftrightarrow
\forall x,y\in A,\; x\leq y\leftrightarrow f(x)\leq' f(y)$.
\end{remark*}

\begin{remark*}[Predicate reading]
$\operatorname{OrderEmbedding}(f,\mathsf{PartiallyOrderedSet}(A,\leq),
\mathsf{PartiallyOrderedSet}(B,\leq'))$ means that $f$ both preserves and
reflects order.
\end{remark*}

\begin{remark*}[Interpretation]
An order embedding realizes the domain poset as an ordered copy inside the
codomain poset.
\end{remark*}

\begin{remark*}[Exposition]
The forward implication says $f$ is order-preserving. The reverse implication
says $f$ is order-reflecting: every order comparison visible between images
already came from a comparison in the domain.
\end{remark*}

\begin{remark*}[Exposition]
If $f : A \to B$ is an order embedding from $(A, \leq)$ to $(B, \leq')$, then
$\leq$ coincides with the induced order $\leq_f$:
\[
x \leq y \;\Longleftrightarrow\; f(x) \leq' f(y) \;\Longleftrightarrow\; x \leq_f y.
\]
Thus an order embedding is exactly a function whose induced order agrees with
the existing order on $A$.
\end{remark*}

\begin{dependencies}
\begin{itemize}
  \item \hyperref[def:ordered-set]{Ordered Set}
  \item \hyperref[def:function]{Function}
\end{itemize}
\end{dependencies}

\begin{propositionbox}{Proposition (Order embeddings are injective)}
\begin{proposition}[Order embeddings are injective]\label{prop:embedding-injective}
Every order embedding $f : (A, \leq) \to (B, \leq')$ is injective.
\hyperref[prf:embedding-injective]{Go to proof.}
\end{proposition}
\end{propositionbox}

\begin{remark*}[Standard quantified statement]
$\operatorname{OrderEmbedding}(f,\mathsf{PartiallyOrderedSet}(A,\leq),
\mathsf{PartiallyOrderedSet}(B,\leq'))\to \operatorname{Injective}(f,A,B)$.
\end{remark*}

\begin{remark*}[Predicate reading]
$\operatorname{Injective}(f,A,B)$ follows from order reflection and
antisymmetry in the domain.
\end{remark*}

\begin{remark*}[Interpretation]
An order embedding cannot collapse two distinct domain elements.
\end{remark*}

\begin{remark*}[Exposition]
If two elements have the same image, then their images are mutually comparable.
Order reflection brings both comparisons back to the domain, and antisymmetry
forces equality.
\end{remark*}

\begin{dependencies}
\begin{itemize}
  \item \hyperref[def:order-embedding]{Order Embedding}
  \item \hyperref[def:injective]{Injective Function}
  \item \hyperref[def:partial-order]{Partial Order}
\end{itemize}
\end{dependencies}

\begin{propositionbox}{Proposition (Order embedding is isomorphism onto image)}
\begin{proposition}[Order embedding is isomorphism onto image]\label{prop:embedding-iso}
If $f : (A, \leq) \to (B, \leq')$ is an order embedding, then $f$ is an
order isomorphism from $(A, \leq)$ to the suborder $(f(A), \leq'_{f(A)})$.
\hyperref[prf:embedding-iso]{Go to proof.}
\end{proposition}
\end{propositionbox}

\begin{remark*}[Standard quantified statement]
$\operatorname{OrderEmbedding}(f,\mathsf{PartiallyOrderedSet}(A,\leq),
\mathsf{PartiallyOrderedSet}(B,\leq'))\to
\operatorname{OrderIsomorphism}(f,\mathsf{PartiallyOrderedSet}(A,\leq),
\mathsf{PartiallyOrderedSet}(f(A),\leq'_{f(A)}))$.
\end{remark*}

\begin{remark*}[Predicate reading]
$\operatorname{OrderIsomorphism}(f,\mathsf{PartiallyOrderedSet}(A,\leq),
\mathsf{PartiallyOrderedSet}(f(A),\leq'_{f(A)}))$ means that the order
embedding becomes bijective after its codomain is restricted to its image.
\end{remark*}

\begin{remark*}[Interpretation]
Every order embedding identifies the domain with an ordered substructure of
the codomain.
\end{remark*}

\begin{remark*}[Exposition]
The result is the order-theoretic analogue of restricting a function to its
actual image: the embedding may not hit all of $B$, but it hits all of $f(A)$,
and on that suborder it is an order isomorphism.
\end{remark*}

\begin{remark*}[Exposition]
The inclusion $\mathbb{N}\hookrightarrow\mathbb{Z}$ with the standard order on
both sets is an order embedding. The map $n\mapsto 2n$ embeds
$(\mathbb{N},\leq)$ into itself, with image the even natural numbers carrying
the inherited order.
\end{remark*}

\begin{remark*}[Exposition]
Induced orders and embeddings are the order-theoretic counterpart of inherited
substructures. In analysis, these ideas appear when restricting an order from
$\mathbb{R}$ to subsets, intervals, sequences, and function spaces, and when
order properties are transferred between ordered structures.
\end{remark*}

\begin{dependencies}
\begin{itemize}
  \item \hyperref[def:order-embedding]{Order Embedding}
  \item \hyperref[def:image-set]{Image of a Set}
  \item \hyperref[def:suborder]{Suborder}
  \item \hyperref[def:partial-order]{Partial Order}
\end{itemize}
\end{dependencies}
